% Chapter Template

\chapter{Background and Related Work}\doublespacing % Main chapter title

\label{Chapter2}

\lhead{Chapter II. \emph{Background and Related Work}}
\setstretch{1.2}

\section{The RISC-V Instruction Set Architecture}

\subsection{Modular Base ISA \& Extensions (RV64C)}

% differences between Risc-V and ARM




There are a myriad of differences between Risc-V and ARM. % https://www.wevolver.com/article/risc-v-vs-arm source
The most glaring one ist that the Instruction Set Architecture (ISA) of RISC-V is open while ARM's ISA is closed. Having a closed ISA means that it is proprietary and controlled by specific companies, i.e. the ARM Holdings company, which issues the licences to companies.
This leads to reliability but development and research is hindered due to the companys control and licencing cost. Open ISA's are open source, allowing developers to freely develop and customize the architercture to fit their needs. % cite this 
This flexibility has lead to widespread adoption by major tech companies like IBM, Nvidia and Western Digital (WD). 

\subsection{Privilege Architecture \& Execution Modes} % (M-Mode vs. S/U-Mode)

In order to provide secure isolation between the system and the user on the software level, The RISC-V ISA defines a hierarchical model which consists of three different hardware privilege modes: User Mode (U-mode), Supervisor Mode (S-Mode) and Machine Mode (M-Mode). 
These modes dictate the scope of instructions that are accessible, the memory regions and control status registers (CSR).

\subsubsection{Privilege Modes}

The highest level of privilege belongs to the machine mode.
It is the only mode which is mandatory for every RISC-V implementation.
Operations executed in this mode have unrestricted access to the hardware, the physical memory Map and core-level configuration registers.
Since M-Mode code interacts with the microarchitecture without the need for memory virtualisation or access restrictions, low-level firmware, hardware initialization routines and baremetal code runs exclusively on this level.

The supervisor mode (S-Mode) is designed to host modern operating system kernels, such as Linux. It supports paged virtual memory subsystems (such as Sv39) which are managed via the supervisor adress and translation and protection register (satp for short). 
The supervisor mode enforces memory isolation across different user tasks while also restricting the operating system from directly modifying core-critical M-mode parameters.

As the name implies, U-Mode is used for the unprivileged user space applications.
The processes


\section{Microarchitectural Foundations \& Memory Hierarchy}

\subsection{In-Order vs. Out-of-Order Execution Paradigms}

Rocket exhibits the classic 5-stage in-order pipeline: fetch, decode, execute, memory, writeback.
Every instruction is broken into these 5 sequential steps.

\begin{enumerate} % maybe reformat this into subsubsections
  \item \textbf{Fetch} The first step of the pipeline.
  It retrieves the instruction from memory using the adress in the program counter (PC) and increments the PC afterwards.
  Memory can be SRAM or an instruction cache.
  The PC is a register which contains the address that is presented to the instruction memory at the beginning of a clock cycle.
  While the cycle is running, the instruction is read from memory while the next PC is being calculated.
  The PC is being calculated by adding the value of 4 to the current PC and then deciding if that new value or that of a branch/jump calculation is to be used.
  In the original RISC design, every instruction is 4 bytes long.
  \item \textbf{Decode} Decodes the opcode of the instruction and reads the needed values from the operand file.
  If the underlying logic of the decoded instruction cannot be executed, the instruction is not passed down the pipeline and thus the pipeline stalls.
  During a stall, input flip flops won't accept new bits, so no new calculations will take place during this cycle.
  If the decoded instruction is a jump or a branch, the target address is computed while the register file is read.
  In the following cycle, the branch condition is computed.
  \item \textbf{Execute} This stage contains an ALU and a bit shifter.
  It might also contain a multiple cycle multiplier and divider.
    The ALU provides Boolean functions such as XOR and performs the addition and subtraction of integers.
    It also provides status bits which indicate if the result was 0 or if an overflow occcured.
    Instructions can be classed into three different categories based on latency.
    Latency describes the time it takes for an instruction to go from start to finish.
    \begin{itemize}
      \item \textbf{Register-Register-Operations:} Includes add, subtract, compare and logical operations.
        During this step of the pipeline, the two given arguments are passed to a ALU, which spits out the result by the end of the stage.
        These operations have a single-cycle latency.
      \item \textbf{Memory Reference:} Includes all the loads from memory.
        The ALU added the two arguments, which include a register and a constant offset to produce a virtual adress by the end of the cycle.
        These operations have a two-cycle latency.
      \item \textbf{Multi cycle instructions:} Includes dividing and multiplying integers and all floating-point operations.
        These types of operations feed their arguments to the multi-cycle multiply- and divde units.
        The remainder of the pipeline can continue working while the multiply/divide blocks are busy.
        In order to keep the writeback stage simple, these blocks write their results to a seperate set of registers.
    \end{itemize}
  \item \textbf{Memory access:} In this stage, single-cycle latency instructions are simply forwarded to the writeback stage.
    This guarantees that both one and two cycle latency instructions write their results at the same stage of the pipeline so that only a singular port of the register file is used and is always available.
  \item \textbf{Writeback:} During writeback, both single-cycle and two-cycle instructions write their results into the register file.
\end{enumerate}

\subsubsection{Hazards}

When the pipeline produces wrong results we call it a hazard.

\begin{itemize}
  \item \textbf{Structural hazards:} A situation in which two instructions try to use the same ressources at the same time.
    This can be avoided by replicating hardware. % TODO: elaborate
  \item \textbf{Data hazards:} If an instruction attempts to use data before said data is available in the register file, we call that a data hazard.
    \begin{itemize}
      \item \textbf{Operand Forwarding:} If a subsequent instruction needs the result of the one that came before it, the pipeline can stall due to timeconsuming read/write-operations.
        This can be fixed by bypassing the write/read stage and forwarding the result of the previous instruction to the execute stage of the following instruction. % TODO: add an example here to illustrate maybe some graphics
        On actual hardware the control unit of the CPU must implement the logic required to catch these dependencies in order to enact the bypass.
      \item \textbf{Pipeline interlock}
    \end{itemize}
  \item \textbf{Control hazards:} A control hazard occurs when the pipeline can't figure out the next instruction adress which needs to be fetched in the following cycle.
    This happens because the control flow instruction (i.e. onditional branch, direct jump, indirect jump, exception) has not been decoded or executed.
    In a processor with a pipeline, fetching occurs at every clock cycle.
    This can be simply mitigated by waiting until the target adress and branch condition is known. This leads to so called "bubbles" - cycles where nothing is being done (NOP cycles).
    Since processing power is wasted, this is not a desirable solution.
    There are multiple hazard mitigation strategies:
    Static predictions when hardware prediction tables are missing or disabled. \textit{Predict Not-Taken} assumes the sequential fetch (PC + 4).
    \textit{Backward Taken } assumes that a backward branch will be taken.
    A backward branch is one that has a taget adress that is lower than its own adress.
    \textit{Forward Not-Taken} assumes the forward branch will not be taken.
    The Rocket core mitigates control hazards by using a 28-entry wide branch target buffer (BTB),
    a 6-entry wide Return Adress Stack (RAS) and a 512-entry wide branch history table (BHT).
    If the branch target buffer predicts a target during the instruction fetch stage, Rocket redirects the prorgam coutner with minimal delay.
    Branch mispredictions are detected at the memory stage of the pipeline. 
    Upon detecting a branch misprediction, rocket flushes the pipeline, meaning any and all instructions in the decode and execute stage are discarded,
    leading to a fixed 3-cycle misprediction penalty.
    Since BOOM employs a decoupled, out-of-order superscalar pipeline, control hazard prediction is handled differently.
    BOOM fields a two-tier predictor consisting of a fast next-line predictor who generates single-cycle predictions
    and a multiple-stage backing predictor (BTP) who runs a "tagged geometric histroy with loop predictor" algorithm  which can override the next-line predictor at stage F4.
  \end{itemize}

  % =====================================================================================

\subsection{Cache Subsystems, Latency Tolerance \& MSHRs}

\subsubsection{Memory Wall \& The Average Memory Access Time (AMAT) Model}

High performance computing is constrained by what is known as the memory hole program \cite{wulf1995memory}.
This problem describes the growing difference between core processing frequency and off-chip dynamic dram latency.
An arithmetic operation within an execution unit takes only a single cycle to complete but retrieving an operand from off-chip memory causes a latency of tens to hundreds of clock-cycles.
In order to mitigate this problem, modern systems rely on multi-level cache hierarchies, non-blocking cache controllers and latency-tolerant execution paradigms.

We can quantify the effectiveness of a memory system by using the average memory acess time (AMAT) model: 

\begin{equation}
  \text{AMAT} = t_{\text{L1\_hit}} + \text{MR}_{\text{L1}} \times \left( t_{\text{L2\_hit}} + \text{MR}_{\text{L2\_local}} \times t_{\text{DRAM}} \right)
\end{equation}

The letter $t$ denotes the acess latency and MR represents the cache miss rate.
In order to keep AMAT low cache parameters such as line-size, set-associativity and total storage capacity are scaled up.
By increasing associativity and capacity conflict and capacity misses are eliminated,
which makes sure that computationally intensive working sets remain in the fast on-chip static RAM and don't need to be moved to off-chip memory.
Chipyard usually implements secondary caching by using the SiFive inclusive L2 cache generator.

\subsubsection{Blocking versus Non-Blocking Caches}
 Classic scalar processors implement blocking caches.
 If a blocking cache find a miss on a load operation,
 the cache controller halts and further memory operations are blocked.
 This forces the entire pipeline to stall until the line refill completes.
 Subsequently, even istructions that are independent who target lines which are already present in the cache are barred from proceding.

 In order to deal with this issue, non-blocking caches are introduced.\cite{kroft1981lockup}
 These types of caches offer two important microarchitectural capabilities:

 \begin{enumerate} %TODO: write more here
   \item \textbf{Miss-Under-Miss:} The cache maintains multiple concurrent in-flight memory requests to different memory blocks,
     thus directly exploiting memory-level parallelism across the memory bus.
   \item \textbf{Hit-Under-Miss:} The cache controller cobtinues to service cache hits for downstream instructions while an outstanding miss transaction is asynchronously fetched from deeper memory tiers.
 \end{enumerate}

\subsubsection{Miss Status Holding Registers (MSHRs)}

The core foundation of a non-blocking cache is rooted ib the mshr file.
Each miss status holding register is a specialized piece of hardware that tracks entries that store the metadata of an active, outstanding cache refill:
When a load misses in the cache tags and no request for that specific block exists, a free mshr is allocated.
The mshr records the block tag, destination physical register or reorder buffer tag and byte offset before dispatching an acquisition transaction across the system.
This process is called \textit{primary miss processing}.
If subseqeuent instructions reference the same in-flight cache line before the refill returns,
the controller avoids issuing redundant bus rtansactions by instead merging the secondary miss into the existing replay queue of the mshr.
Upon data return, every pending request within the queue are satisfied with a singular transaction. 
This process is called \textit{secondary miss merging}.
If all the allocated mshr's are busy,
the cache controller registers a structural hazard and must subsequently stall or reject further incoming memory requests until an active transaction retires.

\subsubsection{Latency Tolerance in Rocket versus BOOM}

The capability of using non-blocking caches to their full potential is always limited by how processors handle executions. 
The in-order Rocket core uses a baseline blocking L1 data cache.
Even if it is linked to a non-blocking interface, Rockets in-order pipeline logic forces a pipeline interlock the moment a subseqeuent instruction needs the load operand. % what is a a load use hazard?
Long off-chip memory latency means idle pipeline cycles.

%\subsubsection{Latency Tolerance & Memory-Level Parallelism}

%\subsubsection{SiFive Inclusive L2 cache versus Broadcast manager}


% ====================================================================

\subsection{On Chip Interconects (TileLink \& AXi4)}

As system-on-chip \gls{soc} complexity scales upwards the Interconect on the chip becomes a vocal factor when regarding memory bandwidth, latency, and system composability.
When it comes to modern heterogenous architecture where in-order cores like Rocket are combined with out-of-order superscalar pipelins like BOOM, cache controllers and DMA peripherals,
the interconnect needs to guarantee deadlock-free packet transportation while preserving cache coherency.
Chipyard provides two complementary bus protocols: The industry-standard AXI4 interface and the native, cache-coherent TileLink.

\subsubsection{The TileLink Interconnect Protocol}

TileLink is a bus protocol which was specifially designed to support the integration of high performance RISC-V soc. \cite{cook2017tilelink}
TileLink is different from other bus protocols that rely on ad-hoc sideband wiring for multiprocessor synchronization because it is mathematically structured around a directed acyclic graph channel model in order to guarantee deadlock freedom.
TileLink defines five independent, decoupled communication channels from A to E:

\begin{enumerate}
  \item \textbf{Channel A Master $\rightarrow$ Slave:} This channel transmits operation requests like reads (i.e. Get),
    writes such as \textit{PutFullData}, \textit{PutPartialData} and cache permission acquisitions like \textit{AcquireBlock} and \textit{AcquirePerm}.
  \item \textbf{Channel B Slave $\rightarrow$ Master:} Carries coherence probe requests from outer cache managers to interrogate or downgrade lines held within private L1 caches.
  \item \textbf{Channel C Master $\rightarrow$ Slave:} transmits probe acknowledgements like \textit{ProbeAck} and \textit{ProbeAckData} or voluntary writeback evictions (\textit{Release}, \textit{ReleaseData})
\item \textbf{Channel D Slave $\rightarrow$ Master:} Returns read data bursts and transaction acknowledgements.
  \item \textbf{Channel E Master $\rightarrow$ Slave:} Delivers a single-wire final grant acknowledgement from the master, thus resolving serialization races between concurrent requesters without a centralized bus arbitration.
\end{enumerate}

\subsubsection{The AXI4 Protocol}

While TileLink is responsible for internal coherent traffic,
the Advanced eXtensible Interface 4 \gls{axi4} which is part of the ARM AMBA4 specification is the industrial standard for peripheral and off-chip memory integration. \cite{arm2011axi4}
AXI4 seperates read and write operations into five different unidirectional channels:
\textit{Read Adress (AR)} and \textit{Read Data (R)} provide pipelined read transactions and data streaming back towards the master.
\textit{Write Adress (AW)}, \textit{Write Data (W)} and Write Response (B) are responsible for decouplin adress issuance,
payload transmission and final completion handshaking.

AXI4 provides native full-duplex communication due to write- and read data paths posessing dedicated physical signals.
Out-of-order transaction completion is supported via transaction id tags (ARID, AWID, RID, BID)
which enable memory controllers to reorder memory requests in order to maximize DRAM row-buffer hit rates.

\subsubsection{SOC integration and protocol bridging within chipyard}

Inside the chipyard ecosystem processors, memory systems and crossbars are composed by using the diplomacy framework. \cite{amid2020chipyard}
This framework negotiates memory-mapped adress ranges, bus widths and transaction ID allotments at the elaboration phase,
which eliminates manual wiring bugs and bus-width mismatch hazards.

The system hierarchy is organized as follows: % TODO: add a graphic here

\begin{itemize}
  \item Processor tiles (i.e. Rocket and \gls{BOOM}) interface directly with the system bus via TL-C.  %what is tlc?
  \item traffic which targets external memory routes across the memory bus to a coherent L2 cache or a TileLink Broadcast manager (if the L2 is ommitted).
  \item In order to interface with external memory hardware chipyard synthesizes a TileLink to AXI4 protocol bridge.
    This middleman comverts TileLink channel operations into \gls{AXI4} adress/data beats while simultanously mapping transaction IDS,
    which ensures the complete interoperability between open-source processors and proprietary memory controllers.
\end{itemize}

\section{Agile Hardware Construction \& Generator Frameworks}

In the past, hardware development relied on waterfall models, manual register-transfer-level (\gls{rtl}) tinkering and static netlists.
With increasing complexity, these methods became more and more prone to errors and made debugging harder.
In order to mitigate these issues, modern computer architecture design has adopted agile hardware development. \cite{lee2016agile}
This is a development paradigm which was adopted from software engineering.
It emphasizes iterative tapeouts, continous integration, extensive parameterization and programmatic hardware generators.
Instead of having to design fixed, singular hardware instances,
developers construct parameterized software programs that emit optimized hardware netlists which are suited for specific application requirements.

\subsection{Classical RTL (Verilog/VHDL) vs Metaprogrammed Generators (Chisel/FIRRTL)}

For the longest time, digital hardware design has been dominated by \gls{hdl} like Verilog and VHDL.
While these languages are great for describing concrete gate-level and cycle-by-cycle register-transfer operations,
they aren't suited for modern hardware development.

\subsubsection{Limitations of Hardware Description Languages}

\begin{enumerate}
  \item \textbf{Primitive Parameterization:} Verilog restris parameterization to compile-time scalar constants and static structural instantiation loops.
    Complex parameter dependencies, such as dynamically sizing crossbar arbitration logic based on the number of heterogenous cache ports or computing routing tables can't be evaluated natively within Verilog.
  \item \textbf{Lack of abstraction:} \glspl{hdl}lack object-oriented object-oriented inheritance, parametric polymorphism, functional closures and higher-order-functions \gls{hof}. 
    If one wants to reuse structures, such as a parameterized queue or a generic arbiter, he would have to reqrite them or duplicate them by hand for every distinct payload.
  \item \textbf{Fragility of Structural Interconnects:} Building multi-core SoCs in a hardware description language like Verilog requires explicitly wiring tens of thousands of seperate port signals across hierarchical module boundaries.
    Manual changes across every intermediate netlist is needed if a sideband coherence signal is to be added or if an adress bus needs to be widened.
    This means a lot of maintenance overhead and manual work, which is very timeconsuming and prone to errors.
    Engineers rely on fragile helper scripts, often written in Python or Pearl, 
    which lack semantic awareness of the underlying hardware types and need to be maintained on top of the already existing code.
\end{enumerate}

\subsubsection{Chisel}

In order to rise above these architectural limitations, 
scientists at the university UC Berkely have developed  \gls{chisel} (\textbf{C}onstructing \textbf{H}ardware \textbf{i}n a \textbf{S}cala \textbf{E}mbedded \textbf{L}anguage). \cite{bachrach2012chisel}
Chisel is a domain specific language embedded inside the modern statically typed Scala programming language.
It is not a high-level synthesis tool that infers hardware from C code.
Chisel does not simulate hardware. It constructs a hardware elaboration graph in memory.
By executing it in the runtime environment (JVM),
hardware developers gain acess to modern software engineering design methods during the hardware elaboration process, such as: 

\begin{itemize}
  \item \textbf{Object-Oriented Abstraction:} functional blocks inheretr from abstract classes and provide reusable traits (mixins).
    Rocket and \gls{BOOM} tiles share common reusable traits like \textit{HasCoreParameters} and \textit{HasL1HellaCacheParameters}, 
    which enable reusing code across different distinct execution pipelines.
  \item \textbf{Higher-Order-Functions and Metaprogramming:} Combinational and repetitive hardware logic can be expressed in a concive way by using functional programming methods like \textit{map}, \textit{reduce} and pattern matching.
    For example a priority arbiter or an ALU operand decoder that would require a lot of nested lines of code in Verilog can be produced by a single functional pipeline in \gls{Chisel}.
  \item \textbf{Parametric Polymorphism:} Harware modules can be generic regarding data types.
    For example, a single First-In-First-Out (FIFO) buffer class (\texttt{Queue[T <: Data]}) can instatiate queues for single-bit flags,
    64-bit integer values or even complex TileLink coherence message bundles without having to change the underlying control logic.
  \item \textbf{Graph-Based Diplomatic Interconnects:} The diplomacy framework models on-chip communication as a directed graph.
    During an early elaboration pass, connected nodes negotiate bus widths, adress spaces, burst lengths and master/slave ID allocations.
    These nodes automatically generate the required width adapters, FIFOFixers and crossbar arbiters prior to hardware generation.
\end{itemize}


\subsubsection{FIRRTL: The Layer Inbetween}

In the Chisel environment, the executing Scala program elaborates the circuit into the \gls{FIRRTL} (\textbf{F}lexible \textbf{I}ntermediate \textbf{R}epresentation for \textbf{RTL}) \cite{amid2020chipyard}



\subsection{The Chipyard Ecosystem \& Integrated Generators}

\section{Open Source SOC Ecosystems}

\subsection{Survey of Frameworks (Chipyard, PULP/PULPino, OpenTitan)}

\subsection{Comparative Analysis \& Ecosystem Taxonomy}

\section{Simulation, Emulation, and Evaluation Methodologies}

\subsection{Cycle-Exact Simulation (Verilator) vs ISA Emulation (Spike/QEMU)}

%\subsection{FPGA Prototyping & Baremetal Benchmarking Approaches}

\section{Summary and Identified Research Gap}






% source for rocket and boom https://www2.eecs.berkeley.edu/Pubs/TechRpts/2016/Archive/EECS-2016-17.pdf
